\section{Example Application Scenarios}
\label{sec:applications}

This section applies the design to four scenarios: Web archives, manufacturing traceability, digital identity, and reputation for automated agents. Each is a design exercise grounded in the implemented system, and none describes a deployed application. In every scenario, existing operational systems keep their current role, and the architecture adds value where several of its properties meet in the same resource.

\subsection{Decentralized Web Archives}

Web content changes and disappears, and today's archives preserve it as captures held by individual institutions. A conventional site can continue serving pages from its content system, and a projection worker gives selected \glslink{resource}{resources} journal-relative \glslink{path}{paths}. Each capture stores its deterministic representation, interpretation metadata, source path, and projection progress. While the site's current page remains convenient for browsing, the projected resource adds a history that independent archives can retain and verify. No single archive, however, need hold everything. Several may preserve overlapping portions of that history---one holding complete captures, another text and metadata, and a third selected high-value resources---and their \glslink{slice}{slices} remain comparable through shared commitments. Reciprocal \glslink{bridge}{bridges} can then place important source and archive \glslink{head}{heads} into one another's later histories. The resulting inclusion receipts survive the loss of any one service when another custodian retains the required evidence. A client whose \glslink{journal}{journal} can reach such an archive through bridges can select it as the retained responder, which serves only evidence it permanently retains, and verify the returned proof against the archive's authenticated object without contacting the vanished publisher.

Adoption can also proceed gradually, because nothing in the model requires existing systems to change. Static sites, databases, and content-management systems remain operational, and the projection adds an audit path. Native clients can later resolve historical resources directly, follow links into retained versions, or compare captures from several archival viewpoints. Even so, long-term operation demands capture-fidelity policy, privacy review, legal-erasure procedures, malicious-content handling, media renewal, and cryptographic migration~\cite{Chun2009tiered}, and the journal gives those activities a stable structure of record.

\subsection{Production Traceability}

Consider a manufacturer that stores an inspection report for component c17 in its quality-management database. A projection worker encodes the report at \texttt{(*state* components c17 inspection)}, where the path names the continuing production resource and an application-defined \glslink{object}{object} binds this version's bytes, metadata, and interpretation. The manufacturer's journal then commits the \glslink{successor}{successor} state as a new signed head. The projection extends existing traceability models with independently retained history~\cite{Stouffer2022traceability,Hasan2009provenance}. When a buyer requests evidence, the manufacturer slices that head along the inspection path, retaining the report, its definition, and the ancestor digests needed to verify the root. The buyer can \glslink{merge}{merge} and \glslink{pin}{pin} the slice under its procurement policy, while the manufacturer remains free to \glslink{prune}{prune} its own historical copy. If the two maintain a reciprocal bridge, the buyer also includes the manufacturer's head in a later head of its own, and the resulting inclusion proof serves as a receipt that the inspection state preceded the buyer's subsequent history (\cref{sec:bridge-histories}). Protected details travel separately through the exact-object \glslink{route}{route} of \cref{sec:access-paths}, under the manufacturer's own \glslink{authorization}{authorization}. The buyer therefore holds a proof it can present to an auditor or regulator without the manufacturer's further cooperation. A later correction appears as an identifiable successor rather than an overwrite, and the buyer can record acceptance, dispute, or follow-up inspection in its own history. Cryptographic evidence preserves the reports and their ordering, while physical identifiers, inspection practice, and institutional judgment connect that evidence to the component itself.

\subsection{Self-Sovereign Identity Anchors}

Online identities today typically live inside platforms and identity providers, so an account's name, connections, and history remain with the service where they accumulated. A person or organization can instead establish a stable path beneath a selected well-known root and use it as an anchor for a profile hosted by an evolving journal. The profile can contain current endpoints, public keys, service relationships, and links to selected claims, all interpreted by observers from roots they already recognize. Changing providers is then a historical transition, in which the earlier state points toward the successor relationship, the new host publishes its descriptor and key, and observers follow the retained trail according to their own policy. Historical resolution makes that trail concrete. Resolving the anchor path at an earlier index shows the profile as it stood before a change, and its \glslink{digest}{digest} cites that exact version. Key rotation is itself a committed step authorized by the predecessor key, so an observer can check each transition from the head it already retains (\cref{sec:bridge-histories}). While historic name trails developed a close model for changing personal online identifiers~\cite{Maniatis2003hints}, journal history extends the trail to richer state, reciprocal receipts, and application-defined relationships. Anchor operators, witnesses, and communities define how recovery, delegation, impersonation disputes, privacy, and endpoint discovery work. The journal structure leaves these questions to those rules and preserves the evidence they consume: prior keys, authorized transitions, witnessed heads, and competing forks. A person can change hosts while keeping the history through which others recognize the change.

\subsection{Agent Reputation Discovery}

A user can begin from journals and communities already trusted for a task, then follow committed relationships to reporting agents whose reports cite sources, record methods, receive corrections, and accumulate later outcomes. Reputation is therefore computed rather than assigned, because each user's journal derives it as a view under local rules, from inspectable history that can travel across ranking services. Different communities may retain different evidence and value prediction calibration, source transparency, or specialized performance, so their judgments can diverge while still citing common reports and outcomes. Secure per-node logs and deterministic replay offer a related accountability model~\cite{Haeberlen2007peerreview}, and partial journal history makes selected evidence portable across evaluators. Historical \glslink{integrity}{integrity} contributes stable inputs---which claim appeared, which sources accompanied it, which correction followed, and which community endorsed the result---that a chosen journal can combine into a local discovery view and preserve for inspection.

Automated agents make this model both more necessary and more practical. For people, digital \glslink{record}{records} largely reflect a physical world that offers independent checks such as personal acquaintance, institutional standing, and direct observation, whereas for an agent the digital record is the world itself, so the evidence behind its judgments has no external fallback. People typically rely on a platform to condense reputation into a score, since reading the underlying history at scale is impractical, and the method of condensation is rarely portable or open to inspection. Agents face no such constraint. They can read cited reports, sources, and corrections directly, and a reputation rule can be an ordinary \glslink{self-encoding}{self-encoding} object stored in an agent's own journal. Such a rule behaves like any other resource and can be versioned, cited, shared with other agents, and evaluated by any adopter against the same retained evidence. Because the rule is an object, another agent can retrieve it from a historical head and recompute its result from the definition and state in the proof (\cref{sec:shared-computation}) instead of trusting a reported score. Reputation thereby becomes a set of competing, inspectable definitions rather than a single opaque score, and an agent can refine its rule as outcomes accumulate without depending on a ranking service.
